% from ega0-21-fr.tex
\oldpage[0\textsubscript{IV-1}]{168}
\begin{env}[21.6.3]
\label{0.21.6.3-fr}
Supposons que l'on ait un diagramme commutatif de monomorphismes de corps
\[
\xymatrix{
 k'_0\ar[r]&k'\ar[r]&K'\ar[r]&L'\\
 k_0\ar[u]\ar[r]&k\ar[u]\ar[r]&K\ar[u]\ar[r]&L\ar[u]
}
\]
Il résulte alors de (20.6.17, (ii)) que l'on a un homomorphisme canonique
\begin{equation}
\tag{21.6.3.1}
\label{0.21.6.3.1-fr}
\Delta(L/K,k/k_0)\longrightarrow\Delta(L'/K',k'/k'_0)
\end{equation}
avec une propriété de transitivité évidente, de sorte que l'on peut dire que
$\Delta(L/K,k/k_0)$ est un \emph{foncteur} en le quadruplet
$(k_0,k,K,L)$.
\end{env}

\begin{proposition}[21.6.4]
\label{0.21.6.4-fr}
\begin{enumerate}
  \item[\rm (i)] Soient $k\subset k'\subset k''\subset K\subset L$ cinq corps.
  On a une suite exacte d'homomorphismes canoniques
\begin{equation}
\tag{21.6.4.1}
\label{0.21.6.4.1-fr}
0\to\Delta(L/K,k'/k)\to\Delta(L/K,k''/k)\to\Delta(L/K,k''/k')\to0
\end{equation}
et par suite l'égalité
\begin{equation}
\tag{21.6.4.2}
\label{0.21.6.4.2-fr}
d(L/K,k''/k)=d(L/K,k''/k')+d(L/K,k'/k).
\end{equation}
  \item[\rm (ii)] Soient $k_0\subset k\subset K\subset L\subset M$ cinq corps.
  On a une suite exacte d'homomorphismes canoniques
\begin{equation}
\tag{21.6.4.3}
\label{0.21.6.4.3-fr}
0\to\Delta(L/K,k/k_0)\otimes_LM\to\Delta(M/K,k/k_0)
\to\Delta(M/L,k/k_0)\to0
\end{equation}
et par suite l'égalité
\begin{equation}
\tag{21.6.4.4}
\label{0.21.6.4.4-fr}
d(M/K,k/k_0)=d(M/L,k/k_0)+d(L/K,k/k_0).
\end{equation}
\end{enumerate}
\end{proposition}

\begin{proof}
\begin{enumerate}
  \item[(i)] Considérons le diagramme commutatif
\[
\xymatrix{
0\ar[r]&\Upsilon_{K/k'/k}\otimes_KL\ar[r]\ar[d]&
\Upsilon_{K/k''/k}\otimes_KL\ar[r]\ar[d]&
\Upsilon_{K/k''/k'}\otimes_KL\ar[r]\ar[d]&0\\
0\ar[r]&\Upsilon_{L/k'/k}\ar[r]\ar[d]&
\Upsilon_{L/k''/k}\ar[r]\ar[d]&
\Upsilon_{L/k''/k'}\ar[r]\ar[d]&0\\
0\ar[r]&\Delta(L/K,k'/k)\ar[r]&\Delta(L/K,k''/k)\ar[r]&
\Delta(L/K,k''/k')\ar[r]&0
}
\]
où l'on peut considérer les trois lignes comme des \emph{complexes}
$T_1,T_2,T_3$ respectivement~; la suite exacte (21.6.1.1) et la définition
de $\Delta$ (21.6.1.2) montrent que l'on a une suite exacte de complexes
$0\to T_1\to T_2\to T_3\to0$~; appliquons-lui la suite exacte de
